    \expectedproblemsection{Standard Mathematical Problems}

    \begin{expectedproblem}[1988 IMO, Vieta Jumping]
        Let $a$ and $b$ be positive integers s.t. $ab+1$ divides
        $a^2+b^2$. Show that
        \[
            \frac{a^2+b^2}{ab+1}
        \]
        is the square of an integer.
    \end{expectedproblem}

    \begin{solution}
        Let
        \[
            k=\frac{a^2+b^2}{ab+1}\in\N.
        \]
        Suppose, for contradiction, that $k$ is not a perfect square. Among all
        positive integer solutions of
        \[
            a^2+b^2=k(ab+1),
        \]
        choose one for which $a+b$ is minimal, and assume without loss of
        generality that $a\geq b$.

        If $a=b$, then
        \[
            k=\frac{2a^2}{a^2+1}<2,
        \]
        so $k=1$, a contradiction. Hence $a>b$. Regard the equation as a
        quadratic equation in $a$:
        \[
            x^2-kbx+b^2-k=0.
        \]
        One root is $a$, so the other root is
        \[
            c=kb-a\in\Z.
        \]
        We claim that $c>0$. From
        \[
            a(a-kb)=k-b^2,
        \]
        if $a=kb$, then $k=b^2$, contrary to the assumption. If $a>kb$, then
        the left-hand side is at least $a\geq kb\geq k$, whereas the
        right-hand side is less than $k$, which is impossible. Therefore
        $a<kb$, and hence $c>0$.

        Since $k$ is a positive integer that is not a perfect square, we have
        $k\geq2$. Let
        \[
            f(x)=x^2-kbx+b^2-k.
        \]
        Since
        \[
            f(b)=(2-k)b^2-k<0
        \]
        and the two roots are $c$ and $a$ with $a>b$, it follows that
        \[
            0<c<b.
        \]
        By Vi\`ete's formulas, $c$ is the other root, and therefore
        \[
            c^2+b^2=k(cb+1).
        \]
        Thus $(c,b)$ is another positive integer solution with
        \[
            c+b<a+b,
        \]
        contradicting the minimality of $a+b$. Hence $k$ must be a perfect
        square.
    \end{solution}

    \begin{expectedproblem}[Sophomore's Dream]
        Prove the identity
        \[
            \int_0^1 x^{-x}\dd x
            =
            \sum_{n=1}^{\infty}\frac{1}{n^n}.
        \]
    \end{expectedproblem}

    \begin{solution}
        For $0<x\leq1$, we have $-\ln x\geq0$, and hence
        \[
            x^{-x}
            =
            e^{-x\ln x}
            =
            \sum_{n=0}^{\infty}
            \frac{x^n(-\ln x)^n}{n!}.
        \]
        Every term is nonnegative, so Tonelli's theorem from Mathematical
        Concepts permits term-by-term integration:
        \[
            \int_0^1x^{-x}\dd x
            =
            \sum_{n=0}^{\infty}
            \frac{1}{n!}
            \int_0^1x^n(-\ln x)^n\dd x.
        \]
        In the $n$th integral, let $t=-(n+1)\ln x$. Then
        \[
            \int_0^1x^n(-\ln x)^n\dd x
            =
            \frac{1}{(n+1)^{n+1}}
            \int_0^\infty t^ne^{-t}\dd t
            =
            \frac{n!}{(n+1)^{n+1}}.
        \]
        Therefore
        \[
            \int_0^1x^{-x}\dd x
            =
            \sum_{n=0}^{\infty}\frac{1}{(n+1)^{n+1}}
            =
            \boxed{\sum_{n=1}^{\infty}\frac{1}{n^n}}.
        \]
        Numerically, the common value is approximately
        \[
            1.291285997\ldots.
        \]
    \end{solution}

    \begin{expectedproblem}[1982 IMO]
        A function
        \[
            f\colon\Z_{>0}\to\Z_{\geq0}
        \]
        satisfies
        \[
            f(m+n)-f(m)-f(n)\in\{0,1\}
        \]
        for all positive integers $m$ and $n$. Moreover,
        \[
            f(2)=0,
            \qquad
            f(3)>0,
            \qquad
            f(9999)=3333.
        \]
        Determine $f(1982)$.
    \end{expectedproblem}

    \begin{solution}
        Repeatedly applying the given condition shows that, for all positive
        integers $r$ and $t$,
        \[
            r f(t)
            \leq
            f(rt)
            \leq
            r f(t)+r-1.
        \]
        Apply this estimate to
        \[
            N=9999\cdot1982
        \]
        in two different ways. Using $t=1982$ and then $t=9999$ gives
        \[
            9999f(1982)
            \leq
            f(N)
            \leq
            9999f(1982)+9998
        \]
        and
        \[
            1982\cdot3333
            \leq
            f(N)
            \leq
            1982\cdot3333+1981.
        \]
        Consequently,
        \[
            \frac{1982\cdot3333-9998}{9999}
            \leq
            f(1982)
            \leq
            \frac{1982\cdot3333+1981}{9999}.
        \]
        The left endpoint lies strictly between $659$ and $660$, whereas the
        right endpoint lies strictly between $660$ and $661$. Since
        $f(1982)$ is an integer, it follows that
        \[
            \boxed{f(1982)=660}.
        \]
    \end{solution}

    \begin{expectedproblem}[Happy Ending Problem]
        Five points are given in the plane, no three of which are collinear.
        Prove that four of them are the vertices of a convex quadrilateral.
    \end{expectedproblem}

    \begin{solution}
        Using the convex-hull terminology from Mathematical Concepts, consider
        the convex hull of the five points. If the convex hull has at
        least four vertices, any four of its vertices, taken in cyclic order,
        form a convex quadrilateral.

        It remains to consider the case in which the convex hull is a triangle
        with vertices $A,B,C$. Let the remaining two points be $P$ and $Q$;
        both lie in the interior of $\triangle ABC$. The line $PQ$ contains none
        of $A,B,C$, so two of the three vertices lie in the same open half-plane
        determined by $PQ$. Relabel them as $A$ and $B$.

        The triangle $\triangle ABQ$ meets the line $PQ$ only at $Q$, because
        $A$ and $B$ lie strictly on the same side of that line. Hence
        $P\notin\triangle ABQ$. Similarly, $Q\notin\triangle ABP$. Since
        $A$ and $B$ are vertices of the convex hull of all five points, neither
        lies in the triangle determined by the other three points. Consequently,
        all four points $A,B,P,Q$ are vertices of their convex hull, which is a
        convex quadrilateral.
    \end{solution}

    \begin{expectedproblem}[Sylvester--Gallai Theorem]
        Let $S$ be a finite set of points in the plane, not all collinear. Prove
        that there exists a line containing exactly two points of $S$.
    \end{expectedproblem}

    \begin{solution}
        Call a line determined by points of $S$ \emph{ordinary} if it contains
        exactly two points of $S$. Consider all pairs $(P,\ell)$ s.t.
        $P\in S$, the line $\ell$ is determined by at least two points of $S$,
        and $P\notin\ell$. Since $S$ is finite and not collinear, the set of
        positive distances arising in this way is finite and nonempty. Choose
        $(P,\ell)$ so that
        \[
            h=\dist(P,\ell)>0
        \]
        is as small as possible, and let $H$ be the foot of the perpendicular
        from $P$ to $\ell$.

        Suppose, for contradiction, that $\ell$ contains at least three points
        of $S$. We can choose distinct points $A,B\in S\cap\ell$ so that either
        $A=H$, or $A$ and $B$ lie on the same side of $H$ with $A$ closer to
        $H$. In either case,
        \[
            AB<PB.
        \]
        Indeed, if $A=H$, then $AB=HB<PB$; otherwise,
        $AB=HB-HA<HB<PB$.

        Computing the area of $\triangle PAB$ with bases $AB$ and $PB$ gives
        \[
            \frac12 AB\,h
            =
            [PAB]
            =
            \frac12 PB\,\dist(A,PB).
        \]
        Hence
        \[
            \dist(A,PB)
            =
            \frac{AB}{PB}h
            <
            h.
        \]
        However, the line $PB$ is determined by two points of $S$, and
        $A\notin PB$. This contradicts the minimality of $h$.

        Therefore, $\ell$ contains exactly two points of $S$, so $\ell$ is the
        required ordinary line.
    \end{solution}

    \begin{expectedproblem}[Moser's Circle Problem]
        Let $n$ be a positive integer. Place $n$ points on a circle and join
        every pair of points by a chord.
        Assume that no three chords meet at a common point in the interior of
        the circle. Determine, with proof, the maximum number of regions into
        which the chords divide the circle.
    \end{expectedproblem}

    \begin{solution}
        There are
        \[
            \binom{n}{2}
        \]
        chords. Draw them one at a time. If a new chord contains $r$ existing
        interior intersection points, those points divide it into $r+1$
        segments, and each segment splits one existing region into two.
        Therefore the new chord creates $r+1$ new regions.

        It follows that the final number of regions is
        \[
            1
            +
            \left|\{\text{chords}\}\right|
            +
            \left|\{\text{interior intersections}\}\right|.
        \]
        Every interior intersection is determined by four boundary points:
        among the six chords joining those four points, exactly the two
        diagonals of their cyclic quadrilateral intersect in the interior.
        Conversely, every choice of four boundary points produces one such
        intersection. The assumption that no three chords are concurrent
        ensures that all these intersections are distinct. Hence there are
        \[
            \binom{n}{4}
        \]
        interior intersections.

        Thus the maximum number of regions is
        \[
            \boxed{
                1+\binom{n}{2}+\binom{n}{4}
            }.
        \]
        In particular, the first values are
        \[
            1,\ 2,\ 4,\ 8,\ 16,\ 31,\ldots,
        \]
        so the apparent pattern $2^{n-1}$ fails when $n=6$.
    \end{solution}

    \begin{expectedproblem}[1976 IMO]
        Determine, with proof, the largest number that can be written as the
        product of positive integers whose sum is $1976$. The number of
        factors is not prescribed.
    \end{expectedproblem}

    \begin{solution}
        A maximizing product contains no factor $1$, since replacing a pair
        $1,a$ by $a+1$ preserves the sum and increases the product. It also
        contains no factor $k\geq5$, because replacing $k$ by $3$ and $k-3$
        preserves the sum and increases the product:
        \[
            3(k-3)>k.
        \]
        A factor $4$ may be replaced by $2,2$ without changing the product.
        Thus it suffices to use only $2$ and $3$.

        Three factors equal to $2$ should be replaced by two factors equal to
        $3$, since
        \[
            2+2+2=3+3
            \qquad\text{and}\qquad
            2^3<3^2.
        \]
        Hence at most two factors equal to $2$ occur. Since
        \[
            1976=3\cdot658+2,
        \]
        the only optimal remainder pattern is one factor equal to $2$ and
        $658$ factors equal to $3$.

        Therefore, the largest possible product is
        \[
            \boxed{2\cdot3^{658}}.
        \]
    \end{solution}

    \begin{expectedproblem}[2021 Putnam]
        A grasshopper starts at the origin in the coordinate plane and makes a
        sequence of hops. Each hop has length $5$, and after every hop the
        grasshopper is at a point whose coordinates are both integers.
        Determine the smallest number of hops needed to reach
        \[
            (2021,2021).
        \]
    \end{expectedproblem}

    \begin{solution}
        If a hop has displacement vector $\mathbf{d}=(p,q)$, then
        \[
            p^2+q^2=25.
        \]
        Thus the possible absolute coordinate pairs are $(5,0)$ and $(4,3)$,
        in either order. In particular, a single hop can increase the sum of
        the two coordinates by at most
        \[
            4+3=7.
        \]
        Since
        \[
            2021+2021=4042=7\cdot577+3,
        \]
        at least $578$ hops are necessary.

        This bound is attained because
        \[
            (2021,2021)
            =
            288(3,4)+288(4,3)+(5,0)+(0,5).
        \]
        The right-hand side is a sum of
        \[
            288+288+1+1=578
        \]
        legal hop vectors. Therefore, the minimum number of hops is
        \[
            \boxed{578}.
        \]
    \end{solution}

    \begin{expectedproblem}[Buffon's Needle]
        Infinitely many parallel lines are drawn in the plane, with distance
        $d$ between consecutive lines. A needle of length $\ell$, where
        \[
            0<\ell\leq d,
        \]
        is dropped at random. Assume that its midpoint is uniformly
        distributed modulo one strip and, independently, that its direction is
        uniformly distributed. Determine the probability that the needle
        intersects one of the parallel lines.
    \end{expectedproblem}

    \begin{solution}
        Let $X$ be the distance from the midpoint of the needle to the nearest
        parallel line, and let $\Theta$ be the acute angle between the needle
        and the lines. Then
        \[
            0\leq X\leq\frac{d}{2},
            \qquad
            0\leq\Theta\leq\frac{\pi}{2},
        \]
        and $X$ and $\Theta$ are independent and uniform on these intervals.

        For a fixed angle $\theta$, the needle intersects a line exactly when
        its perpendicular half-length reaches the nearest line, that is, when
        \[
            X\leq\frac{\ell}{2}\sin\theta.
        \]
        Since $\ell\leq d$, the right-hand side never exceeds $d/2$.
        Therefore,
        \[
            \begin{aligned}
                \Prob(\text{intersection})
                &=
                \frac{2}{\pi}
                \int_0^{\pi/2}
                \Prob\left(
                    X\leq\frac{\ell}{2}\sin\theta
                \right)\dd\theta\\
                &=
                \frac{2}{\pi}
                \int_0^{\pi/2}
                \frac{\ell\sin\theta}{d}\dd\theta\\
                &=
                \boxed{\frac{2\ell}{\pi d}}.
            \end{aligned}
        \]
    \end{solution}

    \begin{expectedproblem}[2011 IMO, Windmill Problem]
        Let $S$ be a finite set of at least two points in the plane, and assume
        that no three points of $S$ are collinear. A windmill is the following
        process. Start with a line $\ell$ through a point $P\in S$. Rotate
        $\ell$ clockwise about the pivot $P$ until it first meets another point
        $Q\in S$. The point $Q$ then becomes the new pivot, and the line
        continues to rotate clockwise. The process is repeated indefinitely.

        Prove that one can choose the initial point $P$ and the initial line
        $\ell$ so that every point of $S$ is used as a pivot infinitely many
        times.
    \end{expectedproblem}

    \begin{solution}
        Orient the rotating line and call its two open sides the left and right
        sides. When the pivot changes from $T$ to $U$, retain the orientation of
        the line. Immediately after the change, the old pivot $T$ lies on the
        same side on which $U$ lay immediately before the change. Hence the
        ordered numbers of points on the two sides are invariant throughout the
        process, except at the instants when the line contains two points of
        $S$.

        First suppose that
        \[
            |S|=2m+1.
        \]
        For every $T\in S$, there is an oriented line through $T$ with exactly
        $m$ points of $S$ on each side. To see this, rotate a line through $T$
        by an angle $\pi$. The number of points on its left changes by one at
        each encounter with another point of $S$, and its initial and final
        values sum to $2m$; therefore the value $m$ occurs.

        Choose such a balanced line through any point as the initial state. The
        side counts remain $(m,m)$. Fix $T\in S$ and choose a balanced line
        $r$ through $T$. Among all lines parallel to $r$, the line $r$ is the
        unique one that contains a point of $S$ and has $m$ points on each
        side: translating it in either direction changes the balance at the
        first point crossed. Consequently, when the windmill becomes parallel
        to $r$, it must coincide with $r$, so $T$ is a pivot. Thus every point
        is visited during each half-turn of the windmill line.

        Now suppose that
        \[
            |S|=2m.
        \]
        For every $T\in S$, the same half-turn argument gives an oriented line
        through $T$ with $m-1$ points on its left and $m$ points on its right.
        Start the windmill from one such line. Its ordered side counts remain
        $(m-1,m)$. For any prescribed $T$, choose an oriented line $r$ through
        $T$ with these counts. It is the unique oriented translate in its
        direction with these counts, so the windmill must pass through $T$ when
        it reaches that oriented direction. Hence every point is visited during
        each full turn.

        Only finitely many line directions are determined by pairs of points of
        $S$. Therefore the positive angular increments between successive pivot
        changes have a positive minimum. Since the process continues
        indefinitely, the total angle of rotation tends to infinity. The
        preceding half-turn or full-turn argument then applies repeatedly, and
        every point of $S$ is used as a pivot infinitely many times.
    \end{solution}

    \begin{expectedproblem}[Art Gallery Problem]
        A museum has the shape of a simple polygon with $n$ vertices. A guard
        must stand at a vertex of the polygon and can see a point whenever the
        entire line segment joining the guard to that point lies in the
        polygon. Prove that the entire museum can always be guarded by at most
        \[
            \left\lfloor\frac{n}{3}\right\rfloor
        \]
        guards.
    \end{expectedproblem}

    \begin{solution}
        Triangulate the polygon using nonintersecting diagonals between its
        vertices. The graph formed by the polygon sides and the triangulation
        diagonals admits a proper coloring with three colors.

        To prove this, use induction on $n$. The assertion is immediate for a
        triangle. Every triangulated polygon has an ear: a triangle in the
        triangulation with two sides on the boundary. Remove the ear vertex and
        color the remaining triangulated polygon by induction. The two
        neighbors of the removed vertex are adjacent and therefore have
        different colors, so the removed vertex can be assigned the third
        color. This completes the induction.

        Every triangle of the triangulation has one vertex of each color. Of
        the three color classes, one contains at most
        \[
            \left\lfloor\frac{n}{3}\right\rfloor
        \]
        vertices. Place guards at all vertices in that smallest color class.
        Every triangle contains one of these guards, and a guard at a vertex of
        a triangle sees every point of that triangle because a triangle is
        convex. Since the triangles cover the polygon, these guards see the
        entire museum.
    \end{solution}

    \begin{expectedproblem}[Bertrand's Ballot Problem]
        In an election, candidate $A$ receives $p$ votes and candidate $B$
        receives $q$ votes, where
        \[
            p>q.
        \]
        The $p+q$ ballots are counted in a uniformly random order. Determine
        the probability that, after every ballot is counted, the cumulative
        number of votes for $A$ is strictly greater than the cumulative number
        of votes for $B$.
    \end{expectedproblem}

    \begin{solution}
        Represent a vote for $A$ by a horizontal step and a vote for $B$ by a
        vertical step. A counting order is then a lattice path from $(0,0)$ to
        $(p,q)$, and the required condition is that the path remain strictly
        below the diagonal $y=x$ after leaving the origin.

        The first ballot must be for $A$. Remove this first horizontal step and
        translate the remaining path one unit to the left. We obtain a path
        from $(0,0)$ to $(p-1,q)$ that never enters the region $y>x$.
        By the reflection principle, the number of such paths is
        \[
            \binom{p+q-1}{q}
            -
            \binom{p+q-1}{q-1}.
        \]
        Since all
        \[
            \binom{p+q}{q}
        \]
        ballot orders are equally likely, the desired probability is
        \[
            \begin{aligned}
                \frac{
                    \binom{p+q-1}{q}-\binom{p+q-1}{q-1}
                }{
                    \binom{p+q}{q}
                }
                &=
                \frac{p-q}{p+q}.
            \end{aligned}
        \]
        Therefore the probability is
        \[
            \boxed{\frac{p-q}{p+q}}.
        \]
    \end{solution}

    \begin{expectedproblem}[1964 IMO]
        Seventeen people correspond with one another. Every pair discusses
        exactly one of three topics in their correspondence. Prove that there
        are three people s.t. every pair among them discusses the same
        topic.
    \end{expectedproblem}

    \begin{solution}
        Choose one person $P$. Among the $16$ correspondences involving $P$,
        at least
        \[
            \left\lceil\frac{16}{3}\right\rceil=6
        \]
        concern the same topic; call it topic $1$, and let the six other people
        be $P_1,\ldots,P_6$.

        If some pair $P_i,P_j$ also discusses topic $1$, then
        $P,P_i,P_j$ form the required triple. Suppose instead that no
        correspondence among $P_1,\ldots,P_6$ concerns topic $1$. Their mutual
        correspondences therefore use only topics $2$ and $3$.

        It remains to use the elementary fact $R(3,3)\leq6$. Indeed, choose one
        of six vertices. At least three of its five incident edges have the
        same color. If one of the three edges joining their other endpoints has
        that color, there is a monochromatic triangle; otherwise those three
        edges form a triangle of the other color. Thus among
        $P_1,\ldots,P_6$ there is a triangle all of whose edges have topic $2$
        or all of whose edges have topic $3$.
    \end{solution}

    \begin{expectedproblem}[Sperner's Lemma]
        A triangle $ABC$ is triangulated into finitely many smaller triangles.
        Label every vertex by one of $1,2,3$ so that $A,B,C$ receive labels
        $1,2,3$, respectively; vertices on $AB$ receive only $1$ or $2$;
        vertices on $BC$ receive only $2$ or $3$; and vertices on $CA$ receive
        only $3$ or $1$. Interior vertices may receive any label.

        Prove that the number of small triangles whose three vertex labels are
        $1,2,3$ is odd. In particular, at least one such triangle exists.
    \end{expectedproblem}

    \begin{solution}
        Call an edge a $12$-edge if its endpoints have labels $1$ and $2$.
        Along the subdivided side $AB$, the sequence of labels begins with $1$
        and ends with $2$. Hence the number of changes between $1$ and $2$ on
        $AB$ is odd, so $AB$ contains an odd number of $12$-edges. The other
        two sides contain none.

        Count incidences between small triangles and their $12$-edges modulo
        $2$. Every interior $12$-edge is incident with two small triangles and
        contributes $0$ modulo $2$, whereas every boundary $12$-edge contributes
        $1$. Thus the total incidence count is odd.

        A small triangle with labels $1,2,3$ has exactly one $12$-edge. A small
        triangle using only labels $1$ and $2$ has either zero or two such
        edges, and every other small triangle has none. Therefore the parity of
        the incidence count is exactly the parity of the number of fully
        labeled triangles. That number is consequently odd.
    \end{solution}

    \begin{expectedproblem}[Morley's Trisector Theorem]
        In a triangle $ABC$, trisect each interior angle. Let $P$ be the
        intersection of the two trisectors adjacent to $BC$, let $Q$ be the
        intersection of the two trisectors adjacent to $CA$, and let $R$ be the
        intersection of the two trisectors adjacent to $AB$. Prove that
        $\triangle PQR$ is equilateral.
    \end{expectedproblem}

    \begin{solution}
        Write
        \[
            \angle A=3\alpha,
            \qquad
            \angle B=3\beta,
            \qquad
            \angle C=3\gamma.
        \]
        Then
        \[
            \alpha+\beta+\gamma=\frac{\pi}{3}.
        \]
        Let $\mathcal{R}$ be the circumradius of $\triangle ABC$.

        In $\triangle BCP$, the angles at $B$ and $C$ are $\beta$ and
        $\gamma$. The sine rule and the identity
        \[
            \sin(3t)
            =
            4\sin t\sin\left(\frac{\pi}{3}+t\right)
            \sin\left(\frac{\pi}{3}-t\right)
        \]
        give
        \[
            BP
            =
            \frac{BC\sin\gamma}{\sin(\beta+\gamma)}
            =
            8\mathcal{R}\sin\alpha\sin\gamma
            \sin\left(\frac{\pi}{3}+\alpha\right).
        \]
        Similarly, using $AB=2\mathcal{R}\sin(3\gamma)$ in
        $\triangle ABR$,
        \[
            BR
            =
            8\mathcal{R}\sin\alpha\sin\gamma
            \sin\left(\frac{\pi}{3}+\gamma\right).
        \]
        Since $\angle PBR=\beta$, the law of cosines yields
        \[
            \begin{aligned}
                PR^2
                &=
                \bigl(8\mathcal{R}\sin\alpha\sin\gamma\bigr)^2
                \Biggl(
                    \sin^2\left(\frac{\pi}{3}+\alpha\right)
                    +
                    \sin^2\left(\frac{\pi}{3}+\gamma\right)\\
                &\hspace{7em}
                    -2
                    \sin\left(\frac{\pi}{3}+\alpha\right)
                    \sin\left(\frac{\pi}{3}+\gamma\right)
                    \cos\beta
                \Biggr).
            \end{aligned}
        \]
        Because
        \[
            \left(\frac{\pi}{3}+\alpha\right)
            +
            \left(\frac{\pi}{3}+\gamma\right)
            +\beta
            =
            \pi,
        \]
        the expression in parentheses is $\sin^2\beta$. Hence
        \[
            PR
            =
            8\mathcal{R}\sin\alpha\sin\beta\sin\gamma.
        \]
        Cyclically applying the same calculation gives the identical value for
        $PQ$ and $QR$. Therefore $\triangle PQR$ is equilateral.
    \end{solution}

    \begin{expectedproblem}[Steiner Tree for a Square]
        Determine the minimum possible total length of a connected road network
        containing all four vertices of a square of side length $a$. Roads may
        meet at additional points in the interior of the square.
    \end{expectedproblem}

    \begin{solution}
        The standard structure theorem for Euclidean Steiner minimal trees says
        that a shortest network is a tree, every nonterminal branching point
        has degree $3$, and its three incident edges meet at angles of
        $120^\circ$. The possible square topologies have zero, one, or two
        Steiner points. With zero Steiner points, the shortest network is a
        minimum spanning tree of length $3a$. With one Steiner point, one first
        joins three square vertices by their three-terminal Steiner tree and
        then attaches the remaining vertex by an edge of length at least $a$.
        The three-terminal tree for a right isosceles triangle of legs $a$ has
        length
        \[
            a\sqrt{2+\sqrt3},
        \]
        so this case is also longer than the construction below. Thus it
        remains to examine the full topology with two Steiner points, each
        joined to two adjacent vertices and to the other Steiner point.

        Place the square at
        \[
            (0,-a/2),\quad(0,a/2),\quad(a,-a/2),\quad(a,a/2).
        \]
        The two Steiner points lie on the horizontal symmetry axis. Write them
        as
        \[
            S_1=(d,0),
            \qquad
            S_2=(a-d,0).
        \]
        The angle between the two edges from $S_1$ to the left-hand vertices is
        $120^\circ$. Taking their dot product gives
        \[
            \frac{d^2-a^2/4}{d^2+a^2/4}=-\frac12,
        \]
        so
        \[
            d=\frac{a}{2\sqrt3}.
        \]
        Each of the four terminal edges has length
        \[
            \sqrt{d^2+\frac{a^2}{4}}=\frac{a}{\sqrt3},
        \]
        while
        \[
            S_1S_2=a-2d=a-\frac{a}{\sqrt3}.
        \]
        Thus the total length is
        \[
            4\frac{a}{\sqrt3}
            +a-\frac{a}{\sqrt3}
            =
            \boxed{a(1+\sqrt3)}.
        \]
        Since
        \[
            a(1+\sqrt3)
            <
            a\left(1+\sqrt{2+\sqrt3}\right)
            <
            3a,
        \]
        the two-Steiner construction is the global minimum.
    \end{solution}

    \begin{expectedproblem}[1975 IMO]
        Let $A$ be the sum of the decimal digits of $4444^{4444}$, let $B$ be
        the sum of the decimal digits of $A$, and let $C$ be the sum of the
        decimal digits of $B$. Determine $C$.
    \end{expectedproblem}

    \begin{solution}
        Since $4444<10^4$,
        \[
            4444^{4444}<10^{17776}.
        \]
        Hence
        \[
            A\leq9\cdot17776=159984.
        \]
        Thus $A$ has at most six digits, so $B\leq54$. Among the integers from
        $1$ to $54$, the largest possible digit sum is $13$, attained at $49$.
        Therefore
        \[
            1\leq C\leq13.
        \]

        Repeated digit sums preserve a number modulo $9$. Since
        \[
            4444\equiv7\pmod9,
            \qquad
            7^3\equiv1\pmod9,
            \qquad
            4444\equiv1\pmod3,
        \]
        we have
        \[
            C\equiv4444^{4444}\equiv7\pmod9.
        \]
        The only integer between $1$ and $13$ congruent to $7$ modulo $9$ is
        \[
            \boxed{C=7}.
        \]
    \end{solution}

    \begin{expectedproblem}[Golomb's Tromino Theorem]
        Remove an arbitrary unit square from a $2^n\times2^n$ board. Prove that
        the remaining board can be tiled by L-shaped trominoes, each consisting
        of three unit squares.
    \end{expectedproblem}

    \begin{solution}
        We use induction on $n$. For $n=1$, the board is $2\times2$, and the
        three remaining squares form one L-tromino.

        Suppose the statement holds for $n-1$. Divide a $2^n\times2^n$ board
        into four $2^{n-1}\times2^{n-1}$ quadrants. The missing square lies in
        one quadrant. Place one L-tromino at the center so that it covers one
        central square in each of the other three quadrants.

        Each quadrant is now a $2^{n-1}\times2^{n-1}$ board with exactly one
        square missing: the original missing square in one quadrant and a
        central square in each of the other three. By the induction hypothesis,
        every quadrant can be tiled. Together with the central tromino, these
        tilings cover the whole deficient board.
    \end{solution}

    \begin{expectedproblem}
        A rectangle is tiled by finitely many smaller rectangles whose sides
        are parallel to the sides of the large rectangle. Each small rectangle
        has at least one side of integer length. Prove that the large rectangle
        also has at least one side of integer length.
    \end{expectedproblem}

    \begin{solution}
        Place the large rectangle at
        \[
            [0,a]\times[0,b].
        \]
        For any measurable region $E$, define
        \[
            I(E)
            =
            \iint_E
            \sin(2\pi x)\sin(2\pi y)\dd x\dd y.
        \]
        This quantity is additive over a finite dissection.

        Consider a tile
        \[
            [u,u+w]\times[v,v+h].
        \]
        If $w$ is an integer, then
        \[
            \int_u^{u+w}\sin(2\pi x)\dd x=0;
        \]
        if $h$ is an integer, the corresponding integral in $y$ is zero.
        Therefore every tile has invariant $0$, and additivity gives
        \[
            I([0,a]\times[0,b])=0.
        \]
        On the other hand,
        \[
            I([0,a]\times[0,b])
            =
            \frac{(1-\cos(2\pi a))(1-\cos(2\pi b))}{4\pi^2}.
        \]
        Both factors in the numerator are nonnegative. Their product is zero
        only when
        \[
            \cos(2\pi a)=1
            \qquad\text{or}\qquad
            \cos(2\pi b)=1,
        \]
        which means that $a$ or $b$ is an integer.
    \end{solution}

    \begin{expectedproblem}[Hall's Marriage Theorem]
        A standard deck of $52$ cards is divided into $13$ piles of four cards.
        Prove that one can choose exactly one card from each pile so that the
        chosen cards have the thirteen distinct ranks
        \[
            A,2,3,\ldots,10,J,Q,K.
        \]
    \end{expectedproblem}

    \begin{solution}
        Construct a bipartite graph whose left vertices are the $13$ piles and
        whose right vertices are the $13$ ranks. Join a pile to every rank that
        occurs in it. We verify Hall's condition.

        Take any $k$ piles. Together they contain $4k$ cards. If the union of
        their neighboring ranks contained at most $k-1$ ranks, those ranks
        could account for at most
        \[
            4(k-1)
        \]
        cards in the entire deck, since each rank occurs in four suits. This is
        impossible. Thus every $k$ piles collectively contain at least $k$
        distinct ranks.

        Hall's Marriage Theorem therefore gives a matching from all piles to
        distinct ranks. Choosing the matched card from each pile produces the
        required selection.
    \end{solution}

    \begin{expectedproblem}[Zeckendorf's Theorem]
        Define
        \[
            F_1=1,
            \qquad
            F_2=2,
            \qquad
            F_{n+2}=F_{n+1}+F_n.
        \]
        Prove that every positive integer can be written uniquely as a sum of
        nonconsecutive terms of this Fibonacci sequence.
    \end{expectedproblem}

    \begin{solution}
        For existence, apply the greedy algorithm. Given $N>0$, choose the
        largest $F_k\leq N$. Since $N<F_{k+1}=F_k+F_{k-1}$, the remainder
        satisfies
        \[
            0\leq N-F_k<F_{k-1}.
        \]
        Thus the next Fibonacci number selected by the same procedure has index
        at most $k-2$. Repeating this process terminates and produces a sum of
        nonconsecutive Fibonacci numbers.

        For uniqueness, use the identity
        \[
            F_{k-1}+F_{k-3}+F_{k-5}+\cdots<F_k.
        \]
        In fact, the sum on the left is $F_k-1$. Suppose two valid
        representations differ, and let $F_k$ be the largest term appearing in
        exactly one of them. The representation containing $F_k$ exceeds the
        sum of every possible collection of nonconsecutive smaller terms,
        whereas the other representation uses only such smaller terms. They
        therefore cannot have the same value, a contradiction.
    \end{solution}

    \begin{expectedproblem}[1986 Putnam]
        Let $\mathbf{A}$ and $\mathbf{B}$ be $n\times n$ matrices
        over a field, and suppose that
        \[
            \mathbf{A}+\mathbf{B}=\mathbf{A}\mathbf{B}.
        \]
        Prove that
        \[
            \mathbf{A}\mathbf{B}=\mathbf{B}\mathbf{A}.
        \]
    \end{expectedproblem}

    \begin{solution}
        Rearranging the hypothesis gives
        \[
            (\mathbf{A}-\mathbf{I}_n)(\mathbf{B}-\mathbf{I}_n)
            =
            \mathbf{A}\mathbf{B}-\mathbf{A}-\mathbf{B}+\mathbf{I}_n
            =
            \mathbf{I}_n.
        \]
        For square matrices, a one-sided inverse is a two-sided inverse. Hence
        \[
            (\mathbf{B}-\mathbf{I}_n)(\mathbf{A}-\mathbf{I}_n)=\mathbf{I}_n.
        \]
        Expanding this identity gives
        \[
            \mathbf{B}\mathbf{A}-\mathbf{A}-\mathbf{B}=\mathbf{0}.
        \]
        Comparing with the original relation yields
        \[
            \mathbf{B}\mathbf{A}
            =
            \mathbf{A}+\mathbf{B}
            =
            \mathbf{A}\mathbf{B}.
        \]
    \end{solution}

    \begin{expectedproblem}[Erd\H{o}s--Szekeres Monotone Subsequence Theorem]
        Let
        \[
            a_1,a_2,\ldots,a_{n^2+1}
        \]
        be distinct real numbers. Prove that the sequence contains an
        increasing subsequence of length $n+1$ or a decreasing subsequence of
        length $n+1$.
    \end{expectedproblem}

    \begin{solution}
        For each index $i$, let $p_i$ be the length of the longest increasing
        subsequence ending at $a_i$, and let $q_i$ be the length of the longest
        decreasing subsequence ending at $a_i$.

        Suppose that neither required subsequence exists. Then
        \[
            1\leq p_i\leq n,
            \qquad
            1\leq q_i\leq n
        \]
        for every $i$. There are only $n^2$ possible ordered pairs
        $(p_i,q_i)$, so two indices $i<j$ have
        \[
            (p_i,q_i)=(p_j,q_j).
        \]
        Since the terms are distinct, either $a_i<a_j$ or $a_i>a_j$. In the
        first case, an increasing subsequence ending at $a_i$ can be extended by
        $a_j$, so $p_j>p_i$. In the second case, $q_j>q_i$. Both alternatives
        contradict equality of the ordered pairs.
    \end{solution}

    \begin{expectedproblem}[2002 Putnam]
        Given any five distinct points on a sphere, prove that some four of them
        lie in a closed hemisphere.
    \end{expectedproblem}

    \begin{solution}
        Choose two of the points that are not antipodal, and draw a great circle
        through them. The two chosen points lie on the boundary of each of the
        two closed hemispheres determined by that great circle.

        Each of the remaining three points belongs to at least one of the two
        hemispheres. By the pigeonhole principle, one hemisphere contains at
        least two of those three points. Together with the two boundary points,
        it contains at least four of the five points.
    \end{solution}

    \begin{expectedproblem}[Prouhet--Thue--Morse Partition]
        For a positive integer $m$, partition
        \[
            \{0,1,\ldots,2^m-1\}
        \]
        into sets $A$ and $B$ s.t.
        \[
            \sum_{a\in A}a^r
            =
            \sum_{b\in B}b^r
        \]
        for every $r=0,1,\ldots,m-1$.
    \end{expectedproblem}

    \begin{solution}
        Put an integer $j$ in $A$ when the number $s_2(j)$ of $1$'s in its
        binary expansion is even, and put it in $B$ when $s_2(j)$ is odd.
        Consider
        \[
            P(x)
            =
            \sum_{j=0}^{2^m-1}(-1)^{s_2(j)}x^j
            =
            \prod_{k=0}^{m-1}(1-x^{2^k}).
        \]
        Every factor on the right vanishes at $x=1$, so $P$ has a zero of order
        $m$ there.

        Apply the operator
        \[
            x\frac{\dd}{\dd x}
        \]
        exactly $r<m$ times and then set $x=1$. Since the zero has order $m$,
        the result is $0$. On the other hand,
        \[
            \left.
            \left(x\frac{\dd}{\dd x}\right)^r P(x)
            \right|_{x=1}
            =
            \sum_{j=0}^{2^m-1}(-1)^{s_2(j)}j^r.
        \]
        Separating the even and odd binary parities gives the required equal
        power sums.
    \end{solution}

    \begin{expectedproblem}[Cayley's Formula]
        For an integer $n\geq2$, determine the number of trees with vertex set
        \[
            \{1,2,\ldots,n\}.
        \]
    \end{expectedproblem}

    \begin{solution}
        Associate a Pr\"ufer sequence with a labeled tree as follows. Repeatedly
        delete the leaf with the smallest label and record its unique neighbor.
        Stop when two vertices remain. This produces a sequence of length
        $n-2$, every term of which lies in $\{1,\ldots,n\}$.

        Conversely, start with a sequence
        \[
            (p_1,p_2,\ldots,p_{n-2}).
        \]
        At each step, choose the smallest remaining label that does not occur
        in the current sequence, join it to the first term of the sequence, and
        then delete both that label and the first sequence term. After all terms have
        been removed, join the two labels that remain. This reconstructs a
        unique tree and reverses the encoding process.

        Thus labeled trees are in bijection with length-$(n-2)$ sequences over
        an alphabet of $n$ labels. Their number is therefore
        \[
            \boxed{n^{n-2}}.
        \]
    \end{solution}

    \begin{expectedproblem}[Tournament King Theorem]
        In a tournament, every pair of vertices is joined by exactly one
        directed edge. Prove that there is a vertex $A$ s.t. every other
        vertex $B$ is either an out-neighbor of $A$ or an out-neighbor of some
        out-neighbor of $A$.
    \end{expectedproblem}

    \begin{solution}
        Choose a vertex $A$ of maximum outdegree. Let $B$ be any vertex that
        defeats $A$. If no vertex defeated by $A$ defeats $B$, then $B$ defeats
        $A$ and every out-neighbor of $A$. Consequently,
        \[
            \deg^+(B)\geq\deg^+(A)+1,
        \]
        contradicting the maximality of the outdegree of $A$.

        Therefore, for every $B$ that defeats $A$, there is a vertex $C$ s.t. $A$ defeats $C$ and $C$ defeats $B$. Hence $A$ is a king.
    \end{solution}

    \begin{expectedproblem}[Gale--Shapley Stable Marriage Problem]
        Two sets $X$ and $Y$ each contain $n$ people. Every person has a strict
        preference ordering of all members of the other set. Prove that there
        is a perfect matching with no blocking pair, and give an algorithm that
        finds one.
    \end{expectedproblem}

    \begin{solution}
        Use the deferred-acceptance algorithm. While some person $x\in X$ is
        unmatched and has not proposed to everyone, let $x$ propose to the most
        preferred member of $Y$ who has not yet rejected $x$. Each $y\in Y$
        keeps the most preferred proposal received so far and rejects every
        other proposer. A tentative engagement may therefore be replaced by a
        more preferred one.

        Each ordered pair $(x,y)$ is used in at most one proposal, so the
        algorithm terminates after at most $n^2$ proposals. At termination the
        matching is perfect: if some $x$ were unmatched, then $x$ would have
        proposed to every member of $Y$, so every member of $Y$ would be
        holding a proposal; since the two sets have equal size, no member of
        $X$ could then be unmatched.

        Suppose that $(x,y)$ is a blocking pair in the final matching. Since
        $x$ prefers $y$ to the assigned partner, $x$ must have proposed to $y$
        earlier. The person $y$ either rejected $x$ immediately or later
        replaced $x$ by someone preferred to $x$. Since tentative partners of
        $y$ only improve, the final partner is preferred to $x$, contradicting
        that $(x,y)$ blocks the matching.
    \end{solution}

    \begin{expectedproblem}[Buffon's Noodle]
        Parallel lines at distance $D$ apart are drawn in the plane. A
        rectifiable plane curve of total length $L$ is placed with uniformly
        random translation modulo the strips and uniformly random orientation.
        If $N$ is the number of intersections of the curve with the parallel
        lines, determine $\Exp[N]$.
    \end{expectedproblem}

    \begin{solution}
        First consider a line segment of length $\ell$ making an acute angle
        $\theta$ with the parallel lines. Its perpendicular span is
        \[
            \ell\sin\theta.
        \]
        Averaging over a uniformly random offset modulo $D$, the expected number
        of parallel lines crossed by this segment is
        \[
            \frac{\ell\sin\theta}{D}.
        \]
        Averaging over a uniformly random direction gives
        \[
            \frac{2}{\pi}
            \int_0^{\pi/2}
            \frac{\ell\sin\theta}{D}\dd\theta
            =
            \frac{2\ell}{\pi D}.
        \]

        For a polygonal curve, add the intersection counts of its segments.
        Linearity of expectation gives the same formula with $\ell$ replaced by
        the total length. A general rectifiable curve follows by polygonal
        approximation, which is the corresponding special case of the
        Cauchy--Crofton formula. Therefore,
        \[
            \boxed{\Exp[N]=\frac{2L}{\pi D}}.
        \]
    \end{solution}

    \begin{expectedproblem}[1987 IMO]
        Prove that there is no function
        \[
            f\colon\Nzero\to\Nzero
        \]
        s.t.
        \[
            f(f(n))=n+1987
        \]
        for every $n\in\Nzero$.
    \end{expectedproblem}

    \begin{solution}
        Applying $f$ to the given identity in two different ways gives
        \[
            f(n+1987)
            =
            f(f(f(n)))
            =
            f(n)+1987.
        \]
        Thus the residue of $f(n)$ modulo $1987$ depends only on the residue of
        $n$ modulo $1987$.

        For each $r\in\{0,1,\ldots,1986\}$, write
        \[
            f(r)=1987q_r+s_r,
            \qquad
            0\leq s_r<1987.
        \]
        Then
        \[
            r+1987
            =
            f(f(r))
            =
            1987q_r+f(s_r).
        \]
        Reducing modulo $1987$ shows that
        \[
            s_{s_r}=r.
        \]
        Hence $r\mapsto s_r$ is an involution on a set of $1987$ elements.

        This involution has no fixed point. Indeed, if $s_r=r$, then the same
        displayed identity gives
        \[
            1987(q_r+q_r)=1987,
        \]
        so $2q_r=1$, which is impossible. An involution without fixed points
        partitions its underlying set into pairs, but a set of odd cardinality
        cannot be partitioned into pairs. Therefore no such function exists.
    \end{solution}

    \begin{expectedproblem}[1994 IMO]
        Let
        \[
            A=\{a_1,a_2,\ldots,a_m\}\subseteq\{1,2,\ldots,n\}
        \]
        be a set of distinct positive integers. Suppose that whenever
        $x,y\in A$ and $x+y\leq n$, one also has $x+y\in A$. Prove that
        \[
            \frac{a_1+a_2+\cdots+a_m}{m}
            \geq
            \frac{n+1}{2}.
        \]
    \end{expectedproblem}

    \begin{solution}
        Let $d$ be the least element of $A$. Partition $A$ according to residue
        classes modulo $d$. Consider one nonempty class, and let $b$ be its
        least element. Since $d,b\in A$, the closure condition implies that
        \[
            b,
            \quad
            b+d,
            \quad
            b+2d,
            \quad\ldots
        \]
        all belong to $A$ for as long as they do not exceed $n$. Therefore this
        residue class of $A$ is a complete arithmetic progression
        \[
            b,b+d,\ldots,c,
        \]
        where $c$ is the largest integer not exceeding $n$ that is congruent to
        $b$ modulo $d$.

        We have $b\geq d$ and $c>n-d$, so
        \[
            b+c>n.
        \]
        Since $b+c$ is an integer,
        \[
            b+c\geq n+1.
        \]
        The average of the progression is therefore at least $(n+1)/2$.
        Every residue class has average at least this value, so their weighted
        average, which is the average of all elements of $A$, also satisfies
        \[
            \boxed{
                \frac{a_1+a_2+\cdots+a_m}{m}
                \geq
                \frac{n+1}{2}
            }.
        \]
    \end{solution}

    \begin{expectedproblem}[1974 IMO]
        Three cards bear distinct positive integers
        \[
            0<p<q<r.
        \]
        In each round the cards are shuffled and dealt, one to each of three
        players $A,B,C$. Each player receives as many counters as the number on
        the card dealt to that player, and all counters are retained. At least
        two rounds are played.

        At the end, the players have respectively
        \[
            20,\qquad10,\qquad9
        \]
        counters, and in the final round player $B$ received the card bearing
        $r$. Determine who received the card bearing $q$ in the first round.
    \end{expectedproblem}

    \begin{solution}
        Suppose that $t$ rounds were played. Adding the final totals gives
        \[
            t(p+q+r)=20+10+9=39.
        \]
        Since $t\geq2$ and player $C$ received at least one counter per round,
        $t\leq9$. Hence
        \[
            t=3,
            \qquad
            p+q+r=13.
        \]

        Player $A$ received $20$ counters in three rounds, so $r\geq7$. In the
        last round player $B$ received $r$, so the two earlier cards dealt to
        $B$ have total $10-r$. Checking the possibilities with
        $p<q<r$ and $p+q+r=13$ gives the unique choice
        \[
            (p,q,r)=(1,4,8).
        \]
        Indeed, for $r=7,9,$ or $10$, the required sum $10-r$ cannot be formed
        by two cards, while for $r=8$ it forces both earlier cards of $B$ to be
        $p=1$.

        In each of the first two rounds, players $A$ and $C$ therefore received
        $q$ and $r$. To reach a total of $20$, player $A$ must have received
        \[
            r,r,q,
        \]
        over the three rounds. Consequently player $C$ received
        \[
            q,q,p,
        \]
        and in particular received $q$ in the first round. Thus the required
        player is
        \[
            \boxed{C}.
        \]
    \end{solution}

    \begin{expectedproblem}[2004 Putnam]
        A basketball player's cumulative free-throw percentage was below
        $80\%$ at one point in a season and above $80\%$ at a later point.
        Prove that at some intermediate point the cumulative percentage was
        exactly $80\%$.
    \end{expectedproblem}

    \begin{solution}
        After $n$ attempts, let $m$ be the number of successful free throws and
        define
        \[
            D=5m-4n.
        \]
        The cumulative percentage is below, equal to, or above $80\%$ exactly
        when $D$ is negative, zero, or positive, respectively.

        After a successful free throw, $D$ increases by
        \[
            5-4=1,
        \]
        whereas after a missed free throw, $D$ decreases by $4$. Thus the only
        way for the integer $D$ to pass from a negative value to a positive
        value is through $0$: an upward step has size only $1$. Therefore at
        some intermediate point
        \[
            5m-4n=0,
        \]
        and the cumulative percentage was exactly $4/5=80\%$.
    \end{solution}

    \begin{expectedproblem}[2005 Putnam]
        Find a nonzero polynomial $P(x,y)$ in two variables s.t.
        \[
            P(\lfloor a\rfloor,\lfloor2a\rfloor)=0
        \]
        for every real number $a$.
    \end{expectedproblem}

    \begin{solution}
        Write
        \[
            a=\lfloor a\rfloor+\{a\},
            \qquad
            0\leq\{a\}<1.
        \]
        Then
        \[
            \lfloor2a\rfloor-2\lfloor a\rfloor
            =
            \lfloor2\{a\}\rfloor
            \in\{0,1\}.
        \]
        Hence one may take
        \[
            \boxed{
                P(x,y)=(y-2x)(y-2x-1)
            }.
        \]
        This polynomial is nonzero and vanishes for both possible values of
        $y-2x$.
    \end{solution}

    \begin{expectedproblem}[1996 Putnam]
        Let $[n]=\{1,2,\ldots,n\}$. A subset $X\subseteq[n]$ is called
        \emph{selfish} if
        \[
            |X|\in X.
        \]
        A selfish set is minimal if none of its proper subsets is selfish.
        Determine the number of minimal selfish subsets of $[n]$.
    \end{expectedproblem}

    \begin{solution}
        Let $X$ be a minimal selfish set and put $|X|=k$. Then $k\in X$. We
        claim that $X$ contains no integer $j<k$. If it did, one could choose a
        $j$-element subset $Y\subsetneq X$ containing $j$, and then
        \[
            |Y|=j\in Y,
        \]
        so $Y$ would be selfish, contradicting minimality.

        Conversely, any $k$-element set that contains $k$ and whose other
        elements are all greater than $k$ is minimal selfish. Thus, for a fixed
        $k$, there are
        \[
            \binom{n-k}{k-1}
        \]
        choices. Summing over all possible $k$ gives
        \[
            \sum_{k\geq1}\binom{n-k}{k-1}
            =
            \sum_{j\geq0}\binom{n-j-1}{j}
            =
            F_n,
        \]
        where $F_1=F_2=1$ and $F_n=F_{n-1}+F_{n-2}$. Therefore the number of
        minimal selfish sets is
        \[
            \boxed{F_n}.
        \]
    \end{solution}

    \begin{expectedproblem}[2002 AIME I]
        Let $S$ be a set of distinct positive integers s.t.
        \[
            1\in S,
            \qquad
            \max S=2002.
        \]
        Suppose that, for every $x\in S$, the arithmetic mean of the elements
        of $S\setminus\{x\}$ is an integer. Determine the maximum possible
        value of $|S|$.
    \end{expectedproblem}

    \begin{solution}
        Let
        \[
            |S|=k,
            \qquad
            T=\sum_{x\in S}x.
        \]
        The hypothesis says that
        \[
            T-x\equiv0\pmod{k-1}
        \]
        for every $x\in S$. Hence all elements of $S$ are congruent modulo
        $k-1$. Since $1,2002\in S$,
        \[
            k-1\mid2001.
        \]

        Put $d=k-1$. There must be at least $d+1$ distinct positive integers not
        exceeding $2002$ that are congruent to $1$ modulo $d$. The number of
        such integers is $2001/d+1$, so
        \[
            d+1\leq\frac{2001}{d}+1,
            \qquad\text{hence}\qquad
            d^2\leq2001.
        \]
        Since
        \[
            2001=3\cdot23\cdot29,
        \]
        the largest divisor of $2001$ not exceeding $\sqrt{2001}$ is $29$.
        Therefore $k\leq30$.

        Equality is attainable. Take
        \[
            S
            =
            \{1+29j\mid0\leq j\leq28\}\cup\{2002\}.
        \]
        This set has $30$ elements, all congruent to $1$ modulo $29$. Its total
        is also congruent to $1$ modulo $29$, so deleting any one element leaves
        a sum divisible by $29$. Thus the maximum is
        \[
            \boxed{30}.
        \]
    \end{solution}

    \begin{expectedproblem}[1961 \UTokyo{} Entrance Examination]
        In a triangle $ABC$, points $L\in\overline{BC}$,
        $M\in\overline{CA}$, and $N\in\overline{AB}$ satisfy
        \[
            \frac{LC}{BL}
            =
            \frac{MA}{CM}
            =
            \frac{NB}{AN}
            =2.
        \]
        Let
        \[
            P=AL\cap CN,
            \qquad
            Q=AL\cap BM,
            \qquad
            R=BM\cap CN.
        \]
        Determine
        \[
            \frac{[PQR]}{[ABC]}.
        \]
    \end{expectedproblem}

    \begin{solution}
        Area ratios are preserved by affine transformations, so take
        \[
            A=(0,0),
            \qquad
            B=(1,0),
            \qquad
            C=(0,1).
        \]
        The division conditions give
        \[
            L=\left(\frac23,\frac13\right),
            \qquad
            M=\left(0,\frac23\right),
            \qquad
            N=\left(\frac13,0\right).
        \]
        Solving the three pairs of line equations yields
        \[
            P=\left(\frac27,\frac17\right),
            \quad
            Q=\left(\frac47,\frac27\right),
            \quad
            R=\left(\frac17,\frac47\right).
        \]
        Therefore
        \[
            [PQR]
            =
            \frac12
            \left|
                \det
                \begin{pmatrix}
                    2/7 & 1/7\\
                    -1/7 & 3/7
                \end{pmatrix}
            \right|
            =
            \frac{1}{14}.
        \]
        Since $[ABC]=1/2$,
        \[
            \boxed{
                \frac{[PQR]}{[ABC]}=\frac17
            }.
        \]
    \end{solution}

    \begin{expectedproblem}[1995 \KyotoU{} Entrance Examination]
        For a positive integer $n$, let $f(n)$ be the remainder when $n$ is
        divided by $7$, and define
        \[
            g(n)
            =
            3f\left(\sum_{k=1}^{7}k^n\right).
        \]
        Prove that
        \[
            f(n^7)=f(n)
        \]
        for every positive integer $n$. Then determine the maximum possible
        value of $g(n)$.
    \end{expectedproblem}

    \begin{solution}
        Fermat's little theorem gives
        \[
            n^7\equiv n\pmod7
        \]
        for every integer $n$, including multiples of $7$. Hence
        \[
            f(n^7)=f(n).
        \]

        Since $f$ takes values in $\{0,1,\ldots,6\}$, we have
        \[
            g(n)\leq18
        \]
        for every positive integer $n$. This upper bound is attained at $n=6$.
        Indeed, for $1\leq k\leq6$,
        \[
            k^6\equiv1\pmod7,
        \]
        while $7^6\equiv0\pmod7$. Therefore
        \[
            f\left(\sum_{k=1}^{7}k^6\right)=6,
        \]
        and the maximum possible value is
        \[
            \boxed{18}.
        \]
    \end{solution}

    \begin{expectedproblem}[1993 \TokyoTech{} Entrance Examination]
        Let $P(x)$ be a polynomial of degree $n$. Suppose that
        \[
            P(0),P(1),\ldots,P(n)
        \]
        are all integers. Prove that $P(m)$ is an integer for every
        $m\in\Z$.
    \end{expectedproblem}

    \begin{solution}
        We use induction on the degree. The assertion is immediate for constant
        polynomials. Let
        \[
            Q(x)=P(x+1)-P(x).
        \]
        Then $Q$ has degree at most $n-1$, and
        \[
            Q(0),Q(1),\ldots,Q(n-1)
        \]
        are integers. By the induction hypothesis,
        \[
            Q(m)\in\Z
        \]
        for every $m\in\Z$.

        Since $P(0)\in\Z$, for $m>0$ we have
        \[
            P(m)
            =
            P(0)+\sum_{j=0}^{m-1}Q(j)
            \in\Z.
        \]
        For $m<0$,
        \[
            P(m)
            =
            P(0)-\sum_{j=m}^{-1}Q(j)
            \in\Z.
        \]
        Thus $P(m)$ is an integer for every integer $m$.
    \end{solution}

    \begin{expectedproblem}[2009 \Hitotsubashi{} Entrance Examination]
        Let $m,n\geq2$ be integers satisfying
        \[
            m^3+1=n^3+10^3.
        \]
        Determine $m$ and $n$.
    \end{expectedproblem}

    \begin{solution}
        Rearranging and factoring gives
        \[
            (m-n)(m^2+mn+n^2)=999.
        \]
        Put $d=m-n>0$. Then $d\mid999$ and
        \[
            3n^2+3dn+d^2=\frac{999}{d}.
        \]
        Since $n\geq2$, the left-hand side exceeds $d^2$, so $d^3<999$.
        Hence the only possible positive divisors are
        \[
            d=1,3,9.
        \]
        Substitution gives no integer $n\geq2$ for $d=1$, gives $n=9$ for
        $d=3$, and gives only $n=1$ or a negative value for $d=9$. Therefore
        \[
            \boxed{(m,n)=(12,9)}.
        \]
        The common value is the taxicab number
        \[
            12^3+1^3=9^3+10^3=1729.
        \]
    \end{solution}

    \begin{expectedproblem}[2025 \UTokyo{} Entrance Examination]
        For a positive integer $a$, define
        \[
            f_a(x)=x^2+x-a.
        \]
        A square number means the square of a nonnegative integer. First prove
        that if $n$ is a positive integer and $f_a(n)$ is a square number, then
        $n\leq a$.

        Let $N_a$ be the number of positive integers $n$ for which $f_a(n)$ is
        a square number. Prove that
        \[
            N_a=1
            \quad\Longleftrightarrow\quad
            4a+1\text{ is prime}.
        \]
    \end{expectedproblem}

    \begin{solution}
        Suppose that
        \[
            n^2+n-a=m^2
        \]
        for some $m\in\Nzero$. If $n\geq a+1$, then
        \[
            n^2<n^2+n-a<(n+1)^2,
        \]
        which is impossible for a square. Hence $n\leq a$.

        Multiplying the equation by $4$ and completing the square gives
        \[
            (2n+1-2m)(2n+1+2m)=4a+1.
        \]
        Both factors are positive odd integers. Conversely, let
        \[
            uv=4a+1,
            \qquad
            1\leq u\leq v,
        \]
        with $u$ and $v$ odd. Since $uv\equiv1\pmod4$, the two factors are
        congruent modulo $4$, and therefore
        \[
            n=\frac{u+v-2}{4},
            \qquad
            m=\frac{v-u}{4}
        \]
        are nonnegative integers satisfying $f_a(n)=m^2$. This gives a
        bijection between positive factor pairs of $4a+1$ and admissible values
        of $n$.

        The factorization
        \[
            4a+1=1\cdot(4a+1)
        \]
        always gives the solution $n=a$. There is exactly one factor pair if
        and only if $4a+1$ is prime. Hence
        \[
            \boxed{
                N_a=1
                \quad\Longleftrightarrow\quad
                4a+1\text{ is prime}
            }.
        \]
    \end{solution}

    \begin{expectedproblem}[1977 IMO]
        Let
        \[
            f\colon\Z_{>0}\to\Z_{>0}
        \]
        satisfy
        \[
            f(n+1)>f(f(n))
        \]
        for every positive integer $n$. Prove that
        \[
            f(n)=n
        \]
        for every positive integer $n$.
    \end{expectedproblem}

    \begin{solution}
        We first prove, by induction on $m$, that
        \[
            k\geq m
            \quad\Longrightarrow\quad
            f(k)\geq m.
        \]
        The assertion is immediate for $m=1$. Suppose it holds for some $m$,
        and let $k\geq m+1$. Write $k=n+1$, so $n\geq m$. By the induction
        hypothesis,
        \[
            f(n)\geq m,
        \]
        and applying the same hypothesis once more gives
        \[
            f(f(n))\geq m.
        \]
        Hence
        \[
            f(k)=f(n+1)>f(f(n))\geq m,
        \]
        so $f(k)\geq m+1$. The induction is complete. Taking $k=m$ shows that
        \[
            f(n)\geq n
        \]
        for every $n$.

        Applying this inequality with the input $f(n)$ gives
        \[
            f(f(n))\geq f(n).
        \]
        Therefore
        \[
            f(n+1)>f(f(n))\geq f(n),
        \]
        so $f$ is strictly increasing.

        If $f(n)>n$ for some $n$, then $f(n)\geq n+1$. Since $f$ is
        increasing,
        \[
            f(f(n))\geq f(n+1),
        \]
        contradicting the given inequality. Thus $f(n)\leq n$ for every $n$.
        Combined with $f(n)\geq n$, this yields
        \[
            \boxed{f(n)=n}
        \]
        for every positive integer $n$.
    \end{solution}

    \begin{expectedproblem}[1985 IMO]
        Let $n$ and $k$ be relatively prime positive integers with $k<n$. Each
        element of
        \[
            M=\{1,2,\ldots,n-1\}
        \]
        is colored either blue or white. Suppose that
        \begin{problemchoices}
            \item $i$ and $n-i$ have the same color for every $i\in M$;
            \item $i$ and $|i-k|$ have the same color whenever
                  $i\in M$ and $i\neq k$.
        \end{problemchoices}
        Prove that all elements of $M$ have the same color.
    \end{expectedproblem}

    \begin{solution}
        For $1\leq j\leq n-1$, let $r_j$ be the least positive residue of
        $jk$ modulo $n$. Since
        \[
            \gcd(n,k)=1,
        \]
        the numbers
        \[
            r_1,r_2,\ldots,r_{n-1}
        \]
        are precisely the elements of $M$ in some order.

        We show that consecutive terms in this list have the same color. Let
        $1\leq j\leq n-2$. The equality $r_j+k=n$ cannot occur, because it
        would imply
        \[
            (j+1)k\equiv0\pmod n
        \]
        with $1\leq j+1<n$.

        If $r_j+k<n$, then
        \[
            r_{j+1}=r_j+k.
        \]
        Since $r_{j+1}\neq k$, the second coloring condition gives
        \[
            r_{j+1}\sim|r_{j+1}-k|=r_j,
        \]
        where $x\sim y$ means that $x$ and $y$ have the same color.

        If $r_j+k>n$, then
        \[
            r_{j+1}=r_j+k-n.
        \]
        The first condition gives $r_j\sim n-r_j$, and the second gives
        \[
            n-r_j
            \sim
            |n-r_j-k|
            =r_j+k-n
            =r_{j+1}.
        \]
        Thus $r_j\sim r_{j+1}$ in either case. Consequently all the residues
        $r_1,\\ldots,r_{n-1}$, and hence all elements of $M$, have the same
        color.
    \end{solution}

    \begin{expectedproblem}[2001 IMO]
        Twenty-one girls and twenty-one boys took part in a mathematical
        competition. Each contestant solved at most six problems, and every
        girl--boy pair solved at least one problem in common. Prove that some
        problem was solved by at least three girls and at least three boys.
    \end{expectedproblem}

    \begin{solution}
        For each girl--boy pair, choose one problem that both contestants
        solved. Suppose, for contradiction, that no problem was solved by at
        least three girls and at least three boys.

        Call a problem \emph{girl-sparse} if it was solved by at most two
        girls. Fix a boy. The problems solved by this boy must collectively
        cover all $21$ girls. Since he solved at most six problems, at least one
        of those problems was solved by at least three girls. Hence he solved at
        most five girl-sparse problems. Each such problem accounts for at most
        two of his $21$ chosen girl--boy pairs. Therefore, for this boy, at most
        \[
            5\cdot2=10
        \]
        chosen pairs use a girl-sparse problem. Over all boys, there are at
        most
        \[
            21\cdot10=210
        \]
        such pairs.

        By symmetry, at most $210$ chosen pairs use a problem solved by at most
        two boys. Under the contrary assumption, every problem is girl-sparse
        or is solved by at most two boys. Thus all
        \[
            21^2=441
        \]
        chosen pairs would belong to the union of two sets containing at most
        $210$ pairs each. This is impossible because
        \[
            441>210+210.
        \]
        Therefore some problem was solved by at least three girls and at least
        three boys.
    \end{solution}

    \begin{expectedproblem}[2012 Putnam]
        Let
        \[
            d_1,d_2,\ldots,d_{12}
        \]
        be real numbers in the open interval $(1,12)$. Prove that there exist
        distinct indices $i,j,k$ s.t. $d_i,d_j,d_k$ are the side lengths
        of an acute triangle.
    \end{expectedproblem}

    \begin{solution}
        Reorder the numbers so that
        \[
            d_1\leq d_2\leq\cdots\leq d_{12}.
        \]
        Suppose that no three of them are the side lengths of an acute
        triangle. For every $r\geq3$, the three consecutive numbers
        $d_{r-2},d_{r-1},d_r$ then satisfy
        \[
            d_r^2\geq d_{r-1}^2+d_{r-2}^2.
        \]
        Indeed, the reverse strict inequality would imply both the triangle
        inequality and, by the converse of the Pythagorean theorem, that the
        triangle is acute.

        Let $F_1=F_2=1$ and $F_r=F_{r-1}+F_{r-2}$. Since
        \[
            d_1^2>1=F_1,
            \qquad
            d_2^2>1=F_2,
        \]
        induction gives
        \[
            d_r^2>F_r
        \]
        for every $r$. In particular,
        \[
            d_{12}^2>F_{12}=144,
        \]
        so $d_{12}>12$, contradicting $d_{12}\in(1,12)$. Therefore the
        required three indices exist.
    \end{solution}

    \begin{expectedproblem}[2017 Putnam]
        Let $S$ be the smallest set of positive integers satisfying the
        following conditions:
        \begin{problemchoices}
            \item $2\in S$;
            \item if $n^2\in S$, then $n\in S$;
            \item if $n\in S$, then $(n+5)^2\in S$.
        \end{problemchoices}
        Determine all positive integers that do not belong to $S$.
    \end{expectedproblem}

    \begin{solution}
        Let
        \[
            U=\{n\in\Z_{>0}\mid (n\neq1)\land(5\nmid n)\}.
        \]
        The set $U$ contains $2$ and satisfies both closure conditions. Hence,
        by the minimality of $S$,
        \[
            S\subseteq U.
        \]

        We prove the reverse inclusion by showing that every set $T$ satisfying
        the three conditions contains $U$. Combining the second and third
        conditions gives
        \[
            n\in T
            \quad\Longrightarrow\quad
            n+5\in T.
        \]
        Thus $T$ is closed under adding any nonnegative multiple of $5$.

        Starting from $2\in T$, the following chain is valid:
        \[
            2
            \longrightarrow
            49
            \longrightarrow
            54^2
            \longrightarrow
            56^2
            \longrightarrow
            56
            \longrightarrow
            121
            \longrightarrow
            11.
        \]
        Here the first two arrows use the third condition, the arrows from
        $54^2$ to $56^2$ and from $56$ to $121$ use closure under adding
        multiples of $5$, and the square-root arrows use the second condition.
        From $11\in T$, closure under adding multiples of $5$ gives
        \[
            16\in T,
            \qquad
            36\in T.
        \]
        Taking square roots gives $4,6\in T$. Since $4\in T$, adding $5$
        gives $9\in T$, and taking a square root gives $3\in T$. Hence
        \[
            2,3,4,6\in T.
        \]
        Repeatedly adding $5$ now shows that $T$ contains every positive
        integer other than $1$ and the positive multiples of $5$. Thus
        $U\subseteq T$ for every admissible $T$, and in particular
        $U\subseteq S$.

        Therefore the positive integers not in $S$ are exactly
        \[
            \boxed{1\text{ and the positive multiples of }5}.
        \]
    \end{solution}

    \begin{expectedproblem}[Fagnano's Problem]
        Let $ABC$ be an acute triangle. Choose points
        \[
            P\in\overline{BC},
            \qquad
            Q\in\overline{CA},
            \qquad
            R\in\overline{AB}.
        \]
        Determine the triangle $PQR$ whose perimeter
        \[
            PQ+QR+RP
        \]
        is as small as possible.
    \end{expectedproblem}

    \begin{solution}
        Reflect $P$ across the lines $CA$ and $AB$ to points $P_1$ and $P_2$,
        respectively. Since $Q\in CA$ and $R\in AB$,
        \[
            PQ=P_1Q,
            \qquad
            PR=P_2R.
        \]
        Therefore
        \[
            PQ+QR+RP
            =
            P_1Q+QR+RP_2
            \geq
            P_1P_2.
        \]
        Moreover,
        \[
            AP_1=AP_2=AP,
            \qquad
            \angle P_1AP_2=2\angle A,
        \]
        so
        \[
            P_1P_2=2AP\sin A.
        \]
        This quantity is minimized when $AP$ is the altitude from $A$ to
        $BC$.

        Let $D,E,F$ be the feet of the altitudes from $A,B,C$, respectively.
        When $P=D$, the right-angle cyclic quadrilaterals determined by the
        altitude feet show that the reflections $D_1,D_2$ of $D$ across $CA$
        and $AB$ satisfy
        \[
            D_1,E,F,D_2
        \]
        collinear. Hence equality holds in the preceding triangle inequality
        for
        \[
            (P,Q,R)=(D,E,F).
        \]
        Thus the unique minimizing inscribed triangle is the orthic triangle of
        $ABC$.
    \end{solution}

    \begin{expectedproblem}[Fermat--Torricelli Problem]
        Given a triangle $ABC$, determine the point $P$ in the plane for which
        \[
            PA+PB+PC
        \]
        is as small as possible.
    \end{expectedproblem}

    \begin{solution}
        First suppose that every angle of $ABC$ is less than $120^\circ$. The
        standard construction with outward equilateral triangles produces a
        unique interior point $P$ satisfying
        \[
            \angle APB
            =
            \angle BPC
            =
            \angle CPA
            =120^\circ.
        \]
        Let $\mathbf{u}_A,\mathbf{u}_B,\mathbf{u}_C$ be the unit vectors from
        $P$ toward $A,B,C$, respectively. Since these vectors are separated by
        angles of $120^\circ$,
        \[
            \mathbf{u}_A+\mathbf{u}_B+\mathbf{u}_C=\mathbf{0}.
        \]
        For any point $Q$, let $\mathbf{q}$ be the vector from $P$ to $Q$.
        The Cauchy--Schwarz inequality gives
        \[
            QA\geq PA-\mathbf{u}_A\cdot\mathbf{q},
        \]
        and similarly for $B$ and $C$. Adding the three inequalities yields
        \[
            QA+QB+QC\geq PA+PB+PC.
        \]
        Thus this interior point is the minimizer.

        Now suppose, without loss of generality, that
        \[
            \angle BAC\geq120^\circ.
        \]
        Let $\mathbf{u}$ and $\mathbf{v}$ be the unit vectors from $A$ toward
        $B$ and $C$, and let $\mathbf{q}$ be the vector from $A$ to an arbitrary
        point $Q$. Then
        \[
            QB\geq AB-\mathbf{u}\cdot\mathbf{q},
            \qquad
            QC\geq AC-\mathbf{v}\cdot\mathbf{q}.
        \]
        Since
        \[
            \|\mathbf{u}+\mathbf{v}\|
            =
            2\cos\left(\frac{\angle BAC}{2}\right)
            \leq1,
        \]
        we also have
        \[
            QA=\|\mathbf{q}\|
            \geq
            (\mathbf{u}+\mathbf{v})\cdot\mathbf{q}.
        \]
        Adding gives
        \[
            QA+QB+QC\geq AB+AC,
        \]
        with equality at $Q=A$. Therefore the minimizing point is
        \[
            \boxed{
                \begin{cases}
                    \text{the interior point making three }120^\circ
                    \text{ angles},
                    & \text{if all triangle angles are }<120^\circ,\\
                    \text{the vertex whose angle is }\geq120^\circ,
                    & \text{otherwise}.
                \end{cases}
            }
        \]
    \end{solution}

    \begin{expectedproblem}[Napoleon's Theorem]
        On the three sides $BC,CA,AB$ of an arbitrary triangle $ABC$, construct
        equilateral triangles externally. Let $X,Y,Z$ be the centers of the
        equilateral triangles constructed on $BC,CA,AB$, respectively. Prove
        that $XYZ$ is equilateral.
    \end{expectedproblem}

    \begin{solution}
        Identify the plane with $\C$, and let $a,b,c$ be the complex
        coordinates of $A,B,C$, listed counterclockwise. Put
        \[
            \rho=e^{-i\pi/3}.
        \]
        The third vertices of the outward equilateral triangles on
        $BC,CA,AB$ are
        \[
            p=b+\rho(c-b),
            \qquad
            q=c+\rho(a-c),
            \qquad
            r=a+\rho(b-a),
        \]
        respectively. Since the center of an equilateral triangle is its
        centroid, the complex coordinates $x,y,z$ of $X,Y,Z$ are
        \[
            x=\frac{b+c+p}{3},
            \qquad
            y=\frac{c+a+q}{3},
            \qquad
            z=\frac{a+b+r}{3}.
        \]
        Direct simplification gives
        \[
            z-y=-\rho(y-x).
        \]
        The complex number $-\rho$ has modulus $1$ and argument $2\pi/3$.
        Hence the vector from $Y$ to $Z$ is obtained from the vector from $X$
        to $Y$ by a rotation through $120^\circ$, without changing its length.
        Therefore
        \[
            XY=YZ
            \qquad\text{and}\qquad
            \angle XYZ=60^\circ,
        \]
        so $XYZ$ is equilateral.
    \end{solution}

    \begin{expectedproblem}[Butterfly Theorem]
        Let $M$ be the midpoint of a chord $PQ$ of a circle. Draw two other
        chords $AB$ and $CD$ through $M$, and define
        \[
            X=AD\cap PQ,
            \qquad
            Y=BC\cap PQ.
        \]
        Prove that $M$ is the midpoint of $XY$.
    \end{expectedproblem}

    \begin{solution}
        Choose Cartesian coordinates with
        \[
            M=(0,0)
        \]
        and $PQ$ as the $x$-axis. Since the perpendicular bisector of $PQ$
        passes through the center of the circle, its equation has the form
        \[
            x^2+(y-h)^2=R^2.
        \]
        Excluding only immediate degenerate cases, write the two chords through
        $M$ as
        \[
            AB:\ x=\alpha y,
            \qquad
            CD:\ x=\beta y,
            \qquad
            \alpha\neq\beta.
        \]
        Let $y_A,y_B,y_C,y_D$ be the $y$-coordinates of the corresponding
        points. Substitution into the circle equation and Vi\`ete's formulas
        give
        \[
            \frac{y_A+y_B}{y_Ay_B}
            =
            \frac{2h}{h^2-R^2}
            =
            \frac{y_C+y_D}{y_Cy_D}.
        \]

        Write $X=(x_X,0)$ and $Y=(x_Y,0)$. Collinearity of $A,X,D$ and of
        $B,Y,C$ gives
        \[
            x_X(y_A-y_D)
            =
            (\alpha-\beta)y_Ay_D,
        \]
        \[
            x_Y(y_B-y_C)
            =
            (\alpha-\beta)y_By_C.
        \]
        Hence
        \[
            \frac{x_X+x_Y}{\alpha-\beta}
            =
            \frac{
                y_Ay_B(y_C+y_D)-y_Cy_D(y_A+y_B)
            }{
                (y_A-y_D)(y_B-y_C)
            }
            =0.
        \]
        Thus $x_X=-x_Y$, so $X$ and $Y$ are equidistant from the origin $M$.
        Therefore
        \[
            \boxed{MX=MY}.
        \]
    \end{solution}

    \begin{expectedproblem}
        The number $N$ of typhoons arriving in a given year is modeled by a
        Poisson distribution with parameter $1$. Determine the probability that
        at least one typhoon arrives during the next year.
    \end{expectedproblem}

    \begin{solution}
        Since $N\sim\Poi(1)$,
        \[
            \Prob(N\geq1)
            =
            1-\Prob(N=0)
            =
            1-e^{-1}.
        \]
        Therefore, the required probability is
        \[
            \boxed{1-\frac{1}{e}}.
        \]
    \end{solution}

    \begin{expectedproblem}
        Consider the functions
        \[
            z^2,
            \qquad
            \sin z,
            \qquad
            \tan z,
            \qquad
            \RePart(z),
            \qquad
            e^z.
        \]
        Determine:
        \begin{problemsubparts}
            \item which functions are entire;
            \item which function is meromorphic on $\C$ but not entire;
            \item which function is not holomorphic on any nonempty open
                  subset of $\C$.
        \end{problemsubparts}
    \end{expectedproblem}

    \begin{solution}
        The functions
        \[
            z^2,
            \qquad
            \sin z,
            \qquad
            e^z
        \]
        are entire. The function $\tan z$ is meromorphic but not entire because
        it has poles at
        \[
            \frac{\pi}{2}+k\pi,
            \qquad
            k\in\Z.
        \]
        Finally, $\RePart(z)$ is not complex differentiable at any point: its
        difference quotient equals $1$ along real increments and $0$ along
        purely imaginary increments. Thus the answers are
        \[
            \boxed{
                \begin{aligned}
                    \text{(a)}\;&z^2,\ \sin z,\ e^z,\\
                    \text{(b)}\;&\tan z,\\
                    \text{(c)}\;&\RePart(z).
                \end{aligned}
            }
        \]
    \end{solution}

    \begin{expectedproblem}
        In the usual topology on $\R$, determine whether each statement is true
        or false.
        \begin{problemchoices}
            \item There exists a perfect subset of $\R$ containing no
                  nondegenerate interval.
            \item The set $\Z$ is open.
            \item The set $\Z$ is closed.
            \item There exists a countable subset of $\R$ that is neither
                  open nor closed.
        \end{problemchoices}
    \end{expectedproblem}

    \begin{solution}
        The Cantor set is perfect and contains no nondegenerate interval, so
        (1) is true. The set $\Z$ is not open, because no open interval centered
        at an integer lies inside $\Z$, but it is closed because
        $\R\setminus\Z$ is open. Hence (2) is false and (3) is true. Finally,
        $\Q$ is countable and is neither open nor closed in $\R$, so (4) is
        true. Therefore,
        \[
            \boxed{\text{(1) T,\quad (2) F,\quad (3) T,\quad (4) T}}.
        \]
    \end{solution}

    \begin{expectedproblem}
        Which of the following integers can be the order of a nonabelian simple
        group?
        \begin{problemchoices}
            \item $60$;
            \item $48$;
            \item $30$;
            \item $77$.
        \end{problemchoices}
    \end{expectedproblem}

    \begin{solution}
        The alternating group $A_5$ is simple and has order $60$, so $60$ is
        possible.

        A group of order $77=7\cdot11$ has a unique Sylow $11$-subgroup, since
        \[
            n_{11}\equiv1\pmod{11},
            \qquad
            n_{11}\mid7,
        \]
        so it is not simple. If a group of order $30$ were simple, then it would
        have $6$ Sylow $5$-subgroups and $10$ Sylow $3$-subgroups. These would
        already contain
        \[
            6(5-1)+10(3-1)=44
        \]
        nonidentity elements, impossible in a group with only $29$
        nonidentity elements.

        Finally, for a group of order $48$, the number of Sylow $2$-subgroups
        is $1$ or $3$. The first case gives a normal subgroup. In the second
        case, conjugation gives a homomorphism to $S_3$; it cannot be injective
        because $48>6$, so its kernel is a nontrivial proper normal subgroup.
        Hence the only possible order is
        \[
            \boxed{60}.
        \]
    \end{solution}

    \begin{expectedproblem}
        Let
        \[
            \mathbf{F}(x,y,z)=(z,x,y).
        \]
        A plane through the origin intersects the sphere
        \[
            x^2+y^2+z^2=4
        \]
        in a great circle. Determine the plane for which the positively
        oriented circulation of $\mathbf{F}$ around this circle is as large as
        possible, and find the maximum circulation.
    \end{expectedproblem}

    \begin{solution}
        The curl is constant:
        \[
            \nabla\times\mathbf{F}
            =
            (1,1,1).
        \]
        Let $\mathbf{n}$ be the unit normal that determines the positive
        orientation of the circle. The circle bounds a disk of radius $2$ and
        area $4\pi$. By Stokes' Theorem,
        \[
            \oint_{\partial S}\mathbf{F}\cdot\dd\mathbf{r}
            =
            \iint_S
            (\nabla\times\mathbf{F})\cdot\mathbf{n}\dd S
            =
            4\pi(1,1,1)\cdot\mathbf{n}.
        \]
        Cauchy--Schwarz gives
        \[
            (1,1,1)\cdot\mathbf{n}
            \leq
            \sqrt{3},
        \]
        with equality when
        \[
            \mathbf{n}
            =
            \frac{1}{\sqrt3}(1,1,1).
        \]
        Therefore the required plane and maximum circulation are
        \[
            \boxed{x+y+z=0},
            \qquad
            \boxed{4\pi\sqrt3}.
        \]
    \end{solution}

    \begin{expectedproblem}[Cantor's Diagonal Argument]
        Let $\{0,1\}^{\N}$ be the set of all infinite binary sequences. Prove
        that no function
        \[
            f\colon\N\to\{0,1\}^{\N}
        \]
        is surjective.
    \end{expectedproblem}

    \begin{solution}
        Write
        \[
            f(n)=(a_{n1},a_{n2},a_{n3},\ldots),
            \qquad
            a_{nk}\in\{0,1\}.
        \]
        Define a binary sequence $b=(b_1,b_2,\ldots)$ by
        \[
            b_n=1-a_{nn}.
        \]
        For every $n\in\N$, the sequence $b$ differs from $f(n)$ in its
        $n$th coordinate. Hence
        \[
            b\neq f(n)
            \qquad
            \text{for every }n\in\N.
        \]
        Therefore $b$ is not in the image of $f$, so $f$ is not surjective.
    \end{solution}

    \begin{expectedproblem}
        On $\R^2$, determine which of the following functions is not a metric.
        For $\mathbf{x}=(x_1,x_2)$ and $\mathbf{y}=(y_1,y_2)$, let
        \begin{problemchoices}
            \item $d(\mathbf{x},\mathbf{y})
                  =\sqrt{(x_1-y_1)^2+(x_2-y_2)^2}$;
            \item $d(\mathbf{x},\mathbf{y})
                  =\min\{|x_1-y_1|,|x_2-y_2|\}$;
            \item $d(\mathbf{x},\mathbf{y})
                  =|x_1-y_1|+|x_2-y_2|$;
            \item $d(\mathbf{x},\mathbf{y})=0$ if $\mathbf{x}=\mathbf{y}$ and
                  $d(\mathbf{x},\mathbf{y})=1$ otherwise.
        \end{problemchoices}
    \end{expectedproblem}

    \begin{solution}
        Choice (2) is not a metric. For example,
        \[
            \mathbf{x}=(0,0),
            \qquad
            \mathbf{y}=(0,1)
        \]
        are distinct, but
        \[
            d(\mathbf{x},\mathbf{y})
            =
            \min\{0,1\}
            =
            0.
        \]
        Thus the identity-of-indiscernibles axiom fails. The other three
        functions are, respectively, the Euclidean, taxicab, and discrete
        metrics. Therefore, the answer is
        \[
            \boxed{(2)}.
        \]
    \end{solution}

    \begin{expectedproblem}
        Which of the following sets has cardinality different from the other
        four?
        \begin{problemchoices}
            \item the Cantor set;
            \item the set of all functions $f\colon\Q\to\R$;
            \item the set $C([0,1])$ of all continuous functions
                  $[0,1]\to\R$;
            \item $\R$;
            \item the set of all algebraic numbers.
        \end{problemchoices}
    \end{expectedproblem}

    \begin{solution}
        The Cantor set has cardinality $|\R|=\mathfrak{c}$. Also,
        \[
            |\R^{\Q}|
            =
            \mathfrak{c}^{\aleph_0}
            =
            (2^{\aleph_0})^{\aleph_0}
            =
            \mathfrak{c}.
        \]
        A continuous function on $[0,1]$ is determined by its values on the
        countable dense set $\Q\cap[0,1]$, so
        \[
            |C([0,1])|\leq\mathfrak{c}^{\aleph_0}=\mathfrak{c}.
        \]
        Constant functions give the reverse inequality, hence
        $|C([0,1])|=\mathfrak{c}$. By contrast, the algebraic numbers are
        countable because there are only countably many polynomials in
        $\Z[x]$, each with finitely many roots. Therefore the answer is
        \[
            \boxed{\text{the set of all algebraic numbers}}.
        \]
    \end{solution}

    \begin{expectedproblem}
        Determine which of the following series converge:
        \begin{problemchoices}
            \item $\displaystyle\sum_{n=1}^{\infty}\frac1n$;
            \item $\displaystyle
                  \sum_{\substack{p\text{ prime}\\p+2\text{ prime}}}
                  \left(\frac1p+\frac1{p+2}\right)$;
            \item $\displaystyle\sum_{n=1}^{\infty}\frac1{n^2}$;
            \item $\displaystyle\sum_{n=1}^{\infty}\frac1{2^n}$.
        \end{problemchoices}
    \end{expectedproblem}

    \begin{solution}
        The harmonic series in (1) diverges. Brun's Theorem from Mathematical
        Concepts states that the twin-prime reciprocal sum in (2) converges.
        The series in (3) is a $p$-series with $p=2>1$, and the series in (4)
        is geometric with ratio $1/2$, so both converge. Hence
        \[
            \boxed{\text{(2), (3), and (4)}}.
        \]
    \end{solution}

    \begin{expectedproblem}
        A Menger sponge is obtained by repeatedly replacing each cube by $20$
        smaller cubes, each scaled by a factor of $1/3$. Determine its
        self-similar dimension.
    \end{expectedproblem}

    \begin{solution}
        If $d$ is the self-similar dimension, Mathematical Concepts gives
        \[
            20\left(\frac13\right)^d=1.
        \]
        Therefore
        \[
            3^d=20,
            \qquad
            \boxed{d=\frac{\ln20}{\ln3}}.
        \]
    \end{solution}

    \begin{expectedproblem}
        Start with a tetrahedral polyhedron and cut seven pairwise disjoint
        tunnel-like through-holes through it. Let $V$, $E$, and $F$ denote the
        numbers of vertices, edges, and faces in any polyhedral decomposition
        of the resulting boundary surface. Determine
        \[
            V-E+F.
        \]
    \end{expectedproblem}

    \begin{solution}
        The original tetrahedral boundary is a sphere, which has genus $0$ and
        Euler characteristic $2$. Each disjoint through-hole adds one handle,
        so the resulting closed orientable surface has genus
        \[
            g=7.
        \]
        By the Euler-characteristic formula from Mathematical Concepts,
        \[
            V-E+F=2-2g=2-14.
        \]
        Hence
        \[
            \boxed{-12}.
        \]
    \end{solution}

